\chapter[Sermon 40. The Locusts Are Described by a Marvelous Hypotyposis, the Popish Clergy: and It Is Shown, of What Sort the Antichristian War Shall Be.]{The Locusts Are Described by a Marvelous Hypotyposis, the Popish Clergy: and It Is Shown, of What Sort the Antichristian War Shall Be.}
\label{sermon:040}
\markright{Sermon 40. The Locusts Are Described by a Marvelous Hypotyposis ...}

And the likeness of the locusts was like unto horses prepared for battle, and on their heads were as it were crowns, like unto gold: and their faces were as the faces of lions. And they had hair as the hair of women. And their teeth were as the teeth of lions. And they had breastplates, as it were breastplates of iron. And the sound of their wings was as the sound of chariots when many horses run together in battle. And they had tails like unto scorpions, and there were stings in their tails. And their power was to hurt men for five months. And they had a king over them, which is the angel of the bottomless pit, whose name in the Hebrew tongue is Abaddon, but in the Greek Apollyon.

\index{Antichrist}We have already spoken of the origin and power of the locusts. Nevertheless, lest anyone should be hindered by obscurity, so that he could not know the locusts, and beware (for this is the end of the whole prophecy to understand the mysteries of Antichrist, and beware), now also he describes the locusts with a wonderful figure, and their fight against Christ, and against the doctrine of godliness of all other most perilous.

\index{Papacy}And there is no doubt but that the whole army of the Pope is here described, especially the spirituality as they term it. For the soldiers of the Emperor, kings, and all princes serve him, whom they call secular. But in the Pope’s tents of the spiritual army are Cardinals, Patriarchs, Archbishops, Bishops, Abbots, Prelates, and there is no noble of priests, and religious persons of both sexes. Hereunto appertain many universities, Doctors, and Masters, great champions of the Pope: these are verily those locusts of whom the Lord Jesus speaks here. I know how displeasingly many will take this my exposition. And I would gladly (God is my witness) have spared them: but all the blame is in them, which in words and works betray and declare themselves to be locusts. For except the thing itself cry out that those things are done of them, which by the exposition are now brought to light, I will not desire that credit should be given to me. I speak nothing here in the favor of any man, neither for hatred. Let God himself be judge betwixt us, let the truth itself judge. Certainly all expositors with one consent understand by locusts false teachers.

But let us see the description of the Apostle John by the revelation of Jesus Christ, which does no harm to anyone, which does not slander anyone. And he shows the similarities, that is to say, the likenesses of locusts, by which they may be depicted, and as it were set before our eyes, to be like the things which he brings forth. For unto every part he applies a parable or similitude, whereby he expresses most aptly the disposition and manners of the locusts.

\index{Repentance}First, he says that the locusts are like unto horses prepared for battle. By this parable he signifies many things at once: that the clergy should not only be ambitious and proud (for a horse is an image of pride) but moreover rebellious and bold, and even cruelly fierce, and in their incredulity and in all their errors most obstinate. They are utterly ignorant of repentance. For John seems here to have alluded to these words of Jeremiah: “How does it happen that this people is not turned away from so froward an aversion? They cleave stiffly to deceit; they refuse to return. I have marked and heard, and they spoke not right; there was none that was sorry for his evil, and that would say, ‘What have I done?’ Every one of them did run his course, as it were a horse dislodged into battle; certainly, with this kind of people there is no amendment.” They think rather about how they may allure others into errors with them. He signifies moreover that the clergy shall be warlike and the authors of wars, and shall instigate wars against the saints and true worshippers of God. For they have the secular power, as they call it, ready. For a long time now there have been no wars that have not been raised by this kind of people. Histories bear witness hereof. Yea, and in this our time cardinals and bishops have led armies, and so forth. Finally, it is signified hereby that the clergy shall continually vex and weary with spiritual war also, the true church of Christ. Wherefore, in chapter 11, we shall hear how the beast comes out of the bottomless pit and makes war with the excellent prophets of God. They mix therefore and practice as well spiritual as corporal wars. Last of all, it is signified that the popes’ clergy shall be well fed, fair and well-looking, and given to voluptuousness, lusts, and pleasures of the body. For this kind of people represent not horses that are gaunt or lean, such as go to plow and cart, but such as are well kept and fed even to serve upon in the wars. For behold with me and consider of what sort the clergy is (for the most part), and you will say that they are here set forth in their colors.

\index{Rome}Secondly, upon their heads, he says, as it were crowns like unto gold. Rabanus Maurus in chapter 3 of the first book of the Institution of Clerks calls the shaving of the priests’ crown a kingdom, truly a token of the dignity of a king and priest. For priests and monks or friars boast themselves to be kings and priests, and yet in deed are neither of both. For the true faithful before God are kings and priests, 1 Peter 2. But by the ordaining or shaving of the Pope, they receive nothing either of kingdom or priesthood. John, therefore, says upon their heads, as it were crowns like unto gold, for he says not they were crowns, but like as they were crowns of gold. They were not crowns in deed, neither were they due to them. And yet notwithstanding in the end of the world now they have taken upon them diademes, or miters, and crowns of gold also, and the same most precious. Yet have they done this by no right. In times past, Bishops of Rome wore white miters, in token of purity and sincerity, finally of the knowledge of both Testaments; but none of the Apostles nor apostolic men wore them. Therefore they betray themselves like a rat with their own utterance, the which I suppose to be done by God’s providence, that they might be known, and eschewed of Christ's sheep as crowned wolves.

\index{Daniel (Book of)}Their faces were like the faces of men, not like the faces of locusts. Similarly, in Daniel, the Antichrist is attributed the eyes of a man, signifying industry and policy. These men pretend to great humanity; they are furnished with fair speech. One might think that if humanity were truly elevated, it would be found in them. But they feign these things, with the intent that by creeping thus into people's hearts, they may accomplish their purposes and deceive. In cunning, deception, shrewdness, and practice, as they term it, the Popes’ Legates, Ambassadors, Priests, and Religious persons excel all other wise men of the world. They press into all assemblies of all men, seeking to be made privy to all things. They consider everything as a means of accomplishing their purposes, they feign and dissemble everything, and they can easily outwit and beguile even those who are most astute. Moreover, they are learned, witty, eloquent, and wonderfully crafty in all things. The thing itself speaks and testifies that I write the truth.

\index{Apocalypse (Book of)}\index{Augustine, Saint}And they had here, like the appearance of women: by this likeness he notes their wantonness, idleness, whorish apparel, and effeminate minds. For they are combed and adorned, and very finely apparelled, delighting in women’s jewels, wearing costly garments, especially in the church, where they ought above all to show humility and frugality. Which of the Apostles ever went so decked (or rather disguised) in the temple or outside the temple? The excess and costliness of apparel of Priests and Monks gives no place to the costly array of the Persian Kings. Again, the thing itself speaks. Saint Augustine in a commentary on the seventh of the Apocalypse, in this appearance, says he would understand and show, not only an effeminate or womanly sex, but also either or both sexes. This is what he says, which I leave to be construed and examined by others.

In attributing to them also the teeth of lions, he signifies their cruelty against the poor and faithful followers of Christ. They are most cruel in persecutions, and of blood most thirsty, and are not moved in this by any compassion. They destroy all things with the sword; many devise various tortures. They excel in tyranny Busirides and Phalarides; the thing itself speaks again. For if kings, princes, or magistrates would spare these wretched people, the priests and friars cry out that it is not lawful; finally, they incite the minds of all princes and magistrates against Gospel adherents by prescribing forms of inquisitions and oppressions. To this is added that some of them are hoarders, accumulating with insatiable covetousness and religious robberies, kings’ treasures. Again, some other wasters follow, who spread evil-gotten goods and waste them prodigally in riot, debauchery, whoring, or wars. Therefore, the teeth of lions are rightly attributed to them, in the same way as Amos is said to have attributed to the false prophets. They had also breastplates (thoraces), a defense for the breast, called a breastplate or a vanguard. Others interpret it as sureties, but they cover all the body; breastplates properly cover the breast. By this is signified that their hearts should be obstinate and inflexible. They are stiff-necked and rigidly bound, neither departing one hair’s breadth from their errors, but even assert that the sea itself cannot err; yes, and that the Pope cannot err. For they will not abide to be taught and admonished, but plainly say that the Church of Rome has never erred; therefore, there remains nothing but that you must subscribe to it, or else be condemned as a heretic and suffer death. It is signified moreover that these shall be through another man’s protection most safe. For they have their immunities, they have their privileges, they have the secular power always ready to fight at their request, they have their fraternities, fellowships, leagues, and affinities. What should we say that bishops and abbots are the sons, brethren, and cousins of princes? Whoever therefore touches them, he has touched the apple of the prince’s eye. For even for maintaining them and their state, all men fight as it were for life and lands.

To the Locusts moreover are ascribed wings. For they are exalted above the common state of men, while they are taken and accounted the most fortunate and most excellent in the world, and so forth. Yea, and impudently they boast that herein they are worthier and greater than the Virgin Mary, for that she bore one in her womb, the Son of God, but they can call him daily to the altar? And while they fly, they make such a noise as horses do, drawing warlike chariots, and now ready to invade the ranks of enemies: that is to say, all their doings are most vehement, most warlike, to men horrible, and deadly. Hereunto pertain the clamorous disputations of the Sorbonne and other schools, excommunications, sentences given at Rome, the popes’ bulls and writings, the boastings of decrees, and they are in obstinacy invincible. All these things make a noise together, and thunder terribly.

Furthermore, it is added that through their decrees and counsels they will tear apart or invade. Consequently, Daniel also attributes prosperity to Antichristians, saying, “He shall act, and he will prosper.” And they will invade in such a way that, as we have said before, people will desire to die, believing that there is no deliverance.

I have spoken before, in a previous sermon, of the tails of scorpions and of five months. Their venomous doctrine is noted, which nevertheless at certain times shall be reproved, that godly men may beware of it. And who does not see, who does not feel, how grievous or difficult this fight or battle is, waged by such locusts? Therefore, the Lord’s mouth has rightly joined a curse with the locusts. Men would rightly wish to die, so that they might be delivered from such great dangers. Let us weigh and consider these things at this day, and let us pray that we may overcome and escape the most pestilent poison of Antichrist.

For now also the king of these locusts is brought forth, and is pointed out as it were with the finger of Christ. He sets him out by three titles, that he may be better known. The locusts, he says, have a king over them. This king is not lawfully given to them, but they themselves have chosen him. For who does not know that, through the policy of the spiritual fathers, the Pope has been exempted from the jurisdiction of princes, so as to rule over all spirituality? They acknowledge no other magistrate than the Pope of Rome, and they rail upon secular princes (as they call them), and will not obey them. All of them bind themselves and swear allegiance to the see of Rome, which they seek to preserve and maintain, even if all others perish. The form of swearing is known, made by bishops, abbots, and doctors to the Pope. And if kings and princes should so much as touch with their finger one who is anointed with the bishops’ oil, even if he be a church robber, a murderer, a thief, and a parricide, they are held accursed, and they and their realms are excommunicated. Thus I say, the locusts have the Pope as their king.

\index{Jerome, Saint}It is also called the Angel of the bottomless pit [?]. And immediately in the eleventh chapter, he shall be called the beast which ascends out of the bottomless pit. Christ descended to us from heaven, the Angel of the Testament and great counsel. Whoever disregards him rightly hears the angel of the bottomless pit, that is to say, Antichrist sent by Satan himself from hell. For he is the adversary and enemy of Christ, in whom the Devil dwells corporally, as also thought Saint Jerome, that the Devil should wholly inhabit that great Antichrist.

\index{Zechariah (Book of)}Therefore also a true name, and a true title most fitting is given him. For they lie who greet and call him “most blessed father,” “most holy Pope,” and so forth. Christ sets forth this figure with another style and gives him other titles. His name, he says, was Abaddon in Hebrew, and in Greek Apollyon. He publishes his name in either tongue for no other reason than that in either Testament—whereof the one is written in Hebrew, the other in Greek—this title is attributed to him. Abaddon or Abbaddon, or Apollyon signifies a destroyer. But Daniel in chapters 7:8 and 11, and Zechariah in chapter 11, attribute this virtue and property to Antichrist. Paul calls him the son of perdition, that is, most lost, most damnable, and the greatest author of perdition and damnation, which finally shall be to many an author of slaughter, by sundry wars. For through false doctrine he destroys souls, and with tyranny by fire and sword he wastes the land, and those who refuse to obey him, most cruelly. Let the Popes’ acts be considered, and the practices of spiritual fathers; let them be applied to these oracles of God, and then let a comparison and judgment be made. And this is as it were the key, opening to us the sense of this place, and that it should be expounded of Antichrist, whom Paul called the son of perdition. Habad in Hebrew signifies lost or destroyed, and from that comes Habbedon, perdition or destruction. So in Greek Apoleo and Apollume signifies to lose and destroy; from this is Apollyon. The Lord Jesus shall slay this destroyer with the breath of his mouth, and take him away utterly by his glorious coming.
